Hôm nay 13:18
Văn bản đã dán (1)(1).txt
Tài liệu

\documentclass[12pt,a4paper]{article}

\usepackage[utf8]{vietnam}
\usepackage[utf8]{inputenc}
\usepackage[left=2.5cm,right=2.5cm,top=2.5cm,bottom=2.5cm]{geometry}

\usepackage{amsmath,amssymb,amsfonts,amsthm,mathtools}
\usepackage{stmaryrd}
\usepackage[hidelinks]{hyperref}

\theoremstyle{definition}
\newtheorem{dinhnghia}{Định nghĩa}
\renewcommand{\thedinhnghia}{3.\arabic{dinhnghia}}
\newtheorem*{vidu}{Ví dụ}
\newtheorem*{ghichu}{Ghi chú}

\title{Cú pháp và Ngữ nghĩa Hình thức của Object Model\
\large Chương 3 --- \emph{A Precise Approach to Validating UML Models and OCL Constraints} (Mark Richters)}
\author{}
\date{}

\begin{document}
\maketitle

\section{Ký hiệu chuẩn bị}

Giả sử một bảng chữ cái $A$ và $N$ một tập tên được dựng từ bảng chữ cái đó.

[
N \subseteq A^{+} \quad\text{là tập các tên hữu hạn, khác rỗng, được tạo thành từ các chuỗi ký tự trên } A.
]

Một \emph{signature} $\Sigma=(T,\Omega)$ cho kiểu dữ liệu, với $T$ là tập tên kiểu (Integer, Real, Boolean, String, các kiểu collection, v.v.) và $\Omega$ là tập các phép toán trên các kiểu đó; $\Sigma$ được định nghĩa đầy đủ ở Chương~4 và chỉ được xem là cho trước trong chương này.

\section{Cú pháp hình thức của Object Model}

\subsection{Lớp, thuộc tính và operation}

\begin{dinhnghia}\label{def:class}
[
Class \subseteq N.
]
\end{dinhnghia}

$Class$ là tập hữu hạn tên lớp. Mỗi lớp $c\in Class$ sinh ra một kiểu đối tượng (object type) $t_c\in T$ cùng tên với lớp.

\begin{vidu}
[
Class = {Branch, Car, CarGroup, Check, Customer, Employee, Person, Rental, ServiceDepot}.
]
\end{vidu}

\begin{dinhnghia}[Thuộc tính, trang 34--35]\label{def:att}
[
Att_c = {, a : t_c \to t \mid a \in N,\ t \in T ,}.
]
\end{dinhnghia}

Tên thuộc tính trong cùng một lớp phải phân biệt, tức mỗi tên chỉ được gán đúng một kiểu:
[
\forall t, t' \in T:\ \bigl(a:t_c\to t \in Att_c \ \land\ a:t_c\to t' \in Att_c\bigr) \implies t = t'.
]

\begin{vidu}
\begin{align*}
Att_{Person} = {,& firstname : Person\to String,\ lastname : Person\to String,\
& age : Person\to Integer,\ isMarried : Person\to Boolean,\
& email : Person\to Set(String),}.
\end{align*}
\end{vidu}

\begin{dinhnghia}[Operation, trang 36]\label{def:op}
[
Op_c = {, \omega : t_c\times t_1\times\cdots\times t_n \to t ,}.
]
\end{dinhnghia}

Tham số đầu tiên $t_c$ là receiver (\texttt{self}); một operation có thể có nhiều tham số nhưng chỉ có một kiểu kết quả.

\begin{vidu}
[
Op_{Car} = {description : Car\to String},\qquad
Op_{Branch} = {rentalsForDay : Branch\times String \to Set(Rental)}.
]
\end{vidu}

\subsection{Association, role và multiplicity}

\begin{dinhnghia}[Association, trang 36--37]\label{def:assoc}
[
Assoc\subseteq N,\qquad associates : Assoc\to Class^{+},\qquad associates(as)=\langle c_1,\ldots,c_n\rangle,\ n\ge2.
]
\end{dinhnghia}

$associates$ gán mỗi tên association với danh sách các lớp tham gia; $n$ được gọi là bậc (degree) của association.

\begin{vidu}
[
associates(Employment) = \langle Branch, Employee\rangle,\qquad
associates(Maintenance) = \langle ServiceDepot, Check, Car\rangle.
]
\end{vidu}

\begin{dinhnghia}[Role name, trang 38]\label{def:role}
[
roles : Assoc\to N^{+}, \qquad roles(as) = \langle r_1,\ldots,r_n\rangle.
]
\end{dinhnghia}

Các role name trong cùng một association phải phân biệt.

\begin{vidu}
[
roles(Employment) = \langle employer, employee\rangle.
]
\end{vidu}

\begin{dinhnghia}[Multiplicity, trang 40]\label{def:mult}
[
multiplicities(as) = \langle M_1,\ldots,M_n\rangle, \qquad \varnothing\ne M_i\subseteq\mathbb{N}_0,\quad M_i\ne{0}.
]
\end{dinhnghia}

$M_i$ là tập các cardinality được phép tại đầu thứ $i$ của association; ví dụ $0..1$ ứng với ${0,1}$ và $*$ ứng với $\mathbb{N}_0$.

\begin{vidu}
[
multiplicities(Maintenance) = \langle {0,1},\ \mathbb{N}_0,\ \mathbb{N}_0\rangle.
]
\end{vidu}

\subsection{Generalization và full descriptor}

\begin{dinhnghia}[Generalization hierarchy, trang 41]\label{def:gen}
$\prec\ \subseteq Class\times Class$ là một thứ tự bộ phận (partial order) trên $Class$.
\end{dinhnghia}

Với $c_1\prec c_2$, lớp $c_1$ là lớp con (child) của $c_2$, và $c_2$ là lớp cha (parent) của $c_1$.

\begin{vidu}
[
\prec\ = {(Customer, Person),\ (Employee, Person)}.
]
\end{vidu}

\begin{dinhnghia}[Full descriptor của một lớp, trang 41--43]\label{def:fd}
[
parents(c) = {c' \mid c'\in Class \land c\prec c'},
]
[
Att_c^{} = Att_c\cup\bigcup_{c'\in parents(c)}Att_{c'},
\qquad
Op_c^{} = Op_c\cup\bigcup_{c'\in parents(c)}Op_{c'},
]
[
FD_c = (Att_c^{},\ Op_c^{},\ navends^{*}(c)).
]
\end{dinhnghia}

\begin{ghichu}
Định nghĩa này không nằm trong bản tóm tắt \texttt{OCLformal.md}, nhưng được thêm lại ở đây vì ký hiệu $Att_c^{*}$ (tập thuộc tính đầy đủ, kể cả thuộc tính kế thừa) được dùng thẳng ở Định nghĩa~\ref{def:state} (trạng thái hệ thống) mà không được định nghĩa lại. $FD_c$ gồm toàn bộ thuộc tính, operation và role name mà lớp $c$ có, kể cả những gì kế thừa từ các lớp cha.
\end{ghichu}

\subsection{Cú pháp hoàn chỉnh của Object Model}

\begin{dinhnghia}[Cú pháp của Object Model, trang 43--44]\label{def:model}
[
M = (Class,\ Att_c,\ Op_c,\ Assoc,\ associates,\ roles,\ multiplicities,\ \prec).
]
\end{dinhnghia}

Ở đây $Att_c$ và $Op_c$ đại diện cho toàn bộ họ tập hợp ${Att_c}{c\in Class}$ và ${Op_c}{c\in Class}$, mỗi lớp một tập riêng. Đây là kết quả mà một parser cùng bộ kiểm tra kiểu cần tạo ra sau khi giải quyết tên và kiểm tra các ràng buộc well-formedness.

\section{Miền ngữ nghĩa của Object Model}

Cú pháp ở trên chỉ mô tả \emph{schema}. Phần này ánh xạ schema đó sang các object, link và trạng thái hệ thống cụ thể, để rồi kết thúc bằng ngữ nghĩa của toàn bộ Object Model.

\subsection{Định danh đối tượng}

\begin{dinhnghia}[Định danh đối tượng, trang 44]\label{def:oid}
[
oid(c) = {c_1, c_2, \ldots},
]
[
I_{Class}(c) = \bigcup{, oid(c') \mid c'\in Class \land c' \preceq c ,}.
]
\end{dinhnghia}

$oid(c)$ là tập vô hạn các định danh đối tượng tiềm năng của lớp $c$; $\preceq$ là bao đóng phản xạ của $\prec$. Miền $I_{Class}(c)$ gồm định danh của $c$ và của mọi lớp con, thể hiện tính thay thế (substitutability): $c_1\prec c_2 \implies I(c_1)\subseteq I(c_2)$.

\begin{vidu}
[
I(Branch)={Branch_1,Branch_2,\ldots},\qquad
I(Customer)={cu_1,cu_2,\ldots},\qquad
I(Employee)={e_1,e_2,\ldots},
]
[
I(Person)={p_1,p_2,\ldots}\cup I(Customer)\cup I(Employee).
]
\end{vidu}

\subsection{Link}

\begin{dinhnghia}[Link, trang 46]\label{def:link}
Với $associates(as)=\langle c_1,\ldots,c_n\rangle$,
[
I_{Assoc}(as) = I_{Class}(c_1)\times\cdots\times I_{Class}(c_n).
]
Một link là một phần tử $l_{as}\in I_{Assoc}(as)$.
\end{dinhnghia}

\begin{vidu}
[
I(Assignment)=I(Rental)\times I(Car),\qquad
I(Quality)=I(CarGroup)\times I(CarGroup),
]
[
I(Maintenance)=I(ServiceDepot)\times I(Check)\times I(Car).
]
\end{vidu}

\subsection{Trạng thái hệ thống}

\begin{dinhnghia}[Trạng thái hệ thống, trang 47]\label{def:state}
Một trạng thái hệ thống của model $M$ là
[
\sigma(M) = (\sigma_{Class},\ \sigma_{Att},\ \sigma_{Assoc}),
]
trong đó, với $\pi_i$ là phép chiếu thành phần thứ $i$ và $\bar\pi_i$ là phép chiếu mọi thành phần trừ thành phần thứ $i$:
\begin{enumerate}
\item tập object đang tồn tại là hữu hạn: $\sigma_{Class}(c)\subset oid(c)$;
\item mỗi thuộc tính gán một giá trị cho từng object đang tồn tại: $\sigma_{Att}(a):\sigma_{Class}(c)\to I(t)$ với mỗi $a:t_c\to t\in Att_c^{*}$;
\item tập link hiện tại là hữu hạn: $\sigma_{Assoc}(as)\subset I_{Assoc}(as)$;
\item mỗi tập link phải thỏa multiplicity đã khai báo:
[
\forall i\in{1,\ldots,n},\ \forall l\in\sigma_{Assoc}(as):\quad
\bigl|{,l'\mid l'\in\sigma_{Assoc}(as)\land \bar\pi_i(l')=\bar\pi_i(l),}\bigr|\in\pi_i\bigl(multiplicities(as)\bigr).
]
\end{enumerate}
\end{dinhnghia}

\begin{vidu}
Một snapshot nhỏ của model Car Rental, với một Branch và hai Employee:
[
\sigma_{Class}(Branch)={b_1},\qquad
\sigma_{Class}(Employee)={e_1,e_2},
]
[
\sigma_{Att}(location)(b_1)=\text{`Berlin'},\qquad
\sigma_{Att}(salary)(e_1)=7200,\qquad
\sigma_{Att}(salary)(e_2)=3800,
]
[
\sigma_{Assoc}(Management)={(b_1,e_1)},\qquad
\sigma_{Assoc}(Employment)={(b_1,e_1),(b_1,e_2)}.
]
Mọi tập khác của $\sigma(M)$ không được nêu ở trên đều rỗng, chẳng hạn $\sigma_{Class}(Car)=\varnothing$.
\end{vidu}

\subsection{Ý nghĩa của toàn bộ Object Model}

\begin{dinhnghia}[Ngữ nghĩa của Object Model, trang 49]\label{def:om-semantics}
[
\llbracket M\rrbracket_{OM} = {, \sigma \mid \sigma \text{ là một trạng thái hệ thống hợp lệ của } M ,}.
]
\end{dinhnghia}

Đây là bước chốt của chương: một Object Model không được ánh xạ sang một trạng thái duy nhất, mà sang \emph{tập tất cả các trạng thái hệ thống được phép}. Tập $\llbracket M\rrbracket_{OM}$ này là nền cho ngữ nghĩa OCL và cho quan hệ thỏa mãn invariant ở các chương sau.

\end{document}

hãy làm mịn lại các công thức. Trình bày

công thức class diagram

1.1 Class
1.2 Attribute
1.3 Asssoc: gồm các hàm
association, rolename, kèm theo các hàm tiện ích: participant, navends
cuối cùng là mutil

1.4 kế thừa
1.5 các tính chất của kế thừa
1.6 các WF của class diagram

biểu diễn công thức hàm giống như hình, áp dụng cho các hàm khacs

Đúng. Như vậy tôi hiểu chính xác hơn: Role vẫn có inheritance ở mức cú pháp để cho phép truy cập thuộc tính của Role cha, nhưng thuộc tính kế thừa không được copy xuống Role con. Nó được truy cập bằng cách ủy quyền qua một instance của Role cha.

Đây là khác biệt rất quan trọng với inheritance thông thường của Entity/Group.

Với Entity/Group

Nếu

Manager≺Person

theo inheritance thông thường và có hai object

m
1
	​

,m
2
	​

∈σ
Class
	​

(Manager),

thì email của Person trở thành thuộc tính của từng object:

m
1
	​

.emailv
a
ˋ
m
2
	​

.email

là hai vị trí trạng thái độc lập:

σ
Att
	​

(email)(m
1
	​

)

=σ
Att
	​

(email)(m
2
	​

)

là hoàn toàn có thể.

Nói cách khác, inheritance ở đây là structural inheritance + object inclusion.

Với Role thì khác

Giả sử

Manager≺Person,Manager,Person∈Role.

Ở mức cú pháp, ta vẫn có:

Att
Manager
∗
	​

=Att
Manager
	​

∪Att
Person
	​

.

Do đó biểu thức

manager.email

vẫn hợp lệ về mặt kiểu.

Nhưng tại state, một instance Manager không sở hữu một bản sao mới của email. Thay vào đó nó phải liên kết tới một instance Person.

Ví dụ:

σ
Role
	​

(Manager)={m
1
	​

},
σ
Role
	​

(Person)={p
1
	​

},

và:

(m
1
	​

,p
1
	​

)∈σ
RoleGen
	​

.

Khi đó:

m
1
	​

.email=σ
Att
	​

(email)(p
1
	​

)
	​


chứ không phải:

σ
Att
	​

(email)(m
1
	​

).

Đây chính là ý bạn nói: về bản chất có một association giữa hai role instance, nhưng cú pháp truy cập được làm trong suốt như attribute inheritance.

Tôi nghĩ có thể hình thức hóa nó rất đẹp bằng cách tách accessible attributes khỏi stored attributes.

Ở tầng model vẫn giữ nguyên:

Att
r
∗
	​

=Att
r
	​

∪
r
′
∈parents(r)
⋃
	​

Att
r
′
	​

,r∈Role.
	​


Nghĩa là nếu:

a:t
r
′
	​

→t∈Att
r
∗
	​

,

thì Role r được phép sử dụng biểu thức:

self.a.

Nhưng Att
r
∗
	​

 ở đây chỉ nói về khả năng truy cập, không nói rằng a được lưu vật lý trên instance của r.

Quan hệ instance-level của Role generalization

Ta thêm vào state một quan hệ:

σ
RoleGen
	​

⊆
r
1
	​

,r
2
	​

∈Role
r
1
	​

≺r
2
	​

	​

⋃
	​

σ
Role
	​

(r
1
	​

)×σ
Role
	​

(r
2
	​

).

Với mỗi specialization:

r
1
	​

≺r
2
	​

,

mỗi instance của Role con phải tham chiếu đúng một instance của Role cha:

∀x∈σ
Role
	​

(r
1
	​

):∃!y∈σ
Role
	​

(r
2
	​

):(x,y)∈σ
RoleGen
	​

.
	​

(RG-1)

Lưu ý rằng quan hệ này không cần injective. Tức là nhiều Role con có thể dùng chung một Role cha.

Đây chính là điều bạn cần.

Ví dụ:

Manager≺Person,Employee≺Person.

Một người p
1
	​

 đồng thời hiện thực hóa hai vai trò:

σ
Role
	​

(Person)={p
1
	​

},
σ
Role
	​

(Manager)={m
1
	​

},σ
Role
	​

(Employee)={e
1
	​

},

với:

(m
1
	​

,p
1
	​

)∈σ
RoleGen
	​

,(e
1
	​

,p
1
	​

)∈σ
RoleGen
	​

.

Nếu:

email:t
Person
	​

→String,

thì:

m
1
	​

.email=e
1
	​

.email=p
1
	​

.email=σ
Att
	​

(email)(p
1
	​

).
	​


Đây là điểm khác hoàn toàn với class inheritance.

Ngữ nghĩa truy cập attribute

Ta có thể định nghĩa một hàm tìm instance của Role tổ tiên:

ancestor
σ
	​

:{
σ
Role
	​

(r)×Role
(x,r
′
)
	​

⇀
↦
	​

RoleInstance
y
	​


với:

y∈σ
Role
	​

(r
′
),r⪯r
′
,

và x liên kết tới y thông qua σ
RoleGen
∗
	​

.

Khi attribute

a:t
r
′
	​

→t

được khai báo tại Role tổ tiên r
′
, thì với

x∈σ
Role
	​

(r),r⪯r
′
,

ngữ nghĩa của truy cập thuộc tính là:

[[x.a]]
σ
	​

=σ
Att
	​

(a)(ancestor
σ
	​

(x,r
′
)).
	​

(ROLE-ATT)

Còn nếu a được khai báo trực tiếp tại r:

[[x.a]]
σ
	​

=σ
Att
	​

(a)(x).

Như vậy từ bên ngoài, người viết vẫn chỉ thấy:

x.email.

Họ không phải viết:

x.parentRole.email

hay:

x.person.email.

Association instance bị ẩn đi bởi semantics của attribute access.

Tôi sẽ gọi cơ chế này là delegated inheritance hoặc chính xác hơn:

attribute access by role delegation
	​


vì nó có ba đặc điểm:

Entity/Group inheritance
Role inheritance
	​

Att
c
∗
	​

Att
r
∗
	​

	​

attribute được k
e
ˆ
ˊ
 thừa v
a
ˋ
o ch
ı
ˊ
nh object
attribute nh
ı
ˋ
n th
a
ˆ
ˊ
y ở Role con nhưng gi
a
ˊ
 trị được l
a
ˆ
ˊ
y từ Role cha.
	​


Vì vậy, công thức Att
c
∗
	​

 ở tầng cú pháp thực ra có thể dùng chung cho cả ba loại:

Att
c
∗
	​

=Att
c
	​

∪
c
′
∈parents(c)
⋃
	​

Att
c
′
	​

,c∈Class.
	​


và tương tự:

navends
∗
(c)=navends(c)∪
c
′
∈parents(c)
⋃
	​

navends(c
′
).
	​


Sự khác biệt không nằm ở Full Descriptor nữa. Nó nằm hoàn toàn ở Interpretation/System State:

c∈Entity∪Group
c∈Role
	​

⇒
⇒
	​

inheritance by object/domain inclusion,
inheritance by instance delegation.
	​

	​


Theo tôi đây là mô hình sạch nhất cho đúng ý tưởng của bạn.

tức la ở class, họ có cái hàm Att_c là lấy danh sách thuộc tính của class C và  σAtt(salary)(e1 ) là lấy thuộc tính salary của đối tượng e1, họ lấy dễ dàng vì họ viết rằng, họ biết rằng salary này là thuộc tính độc nhất của đối tương e1, tức là không có một đối tượng nào khác cũng có thuộc tinhs salary nyaf, tức là thuộc tính salary là thuộc tính của lớp cha. Mà một role con khác thay đổi nó, thì role khác cũng bị thay đổi ý.

Đúng về hệ quả bạn muốn đạt được, nhưng cần chỉnh một điểm trong cách hiểu về σ
Att
	​

.

Trong mô hình Class thông thường của Richters, một attribute không phải là “độc nhất cho một object”. Attribute là một hàm dùng chung cho tất cả object của lớp khai báo nó. Với

salary:t
Employee
	​

→Real,

ta có ở một state:

σ
Att
	​

(salary):σ
Class
	​

(Employee)→I(Real).

Vì vậy nếu

σ
Class
	​

(Employee)={e
1
	​

,e
2
	​

},

thì cùng một attribute salary được áp dụng lên hai object:

σ
Att
	​

(salary)(e
1
	​

)=7200,σ
Att
	​

(salary)(e
2
	​

)=3800.

Tức là ký hiệu attribute salary là chung, nhưng mỗi object có một giá trị riêng. Đây chính là cách Definition 3.12 định nghĩa σ
Att
	​

.

Điểm quan trọng với inheritance của Class là: nếu Employee kế thừa Person và

email:t
Person
	​

→String,

thì một object e
1
	​

:Employee không cần một object Person riêng. Chính e
1
	​

 được dùng làm đối số:

σ
Att
	​

(email)(e
1
	​

).

Bởi vì Employee object đồng thời nằm trong miền của Person:

I(Employee)⊆I(Person).

Richters cũng nói rõ các giá trị thuộc tính của Employee, ngoài salary, được kế thừa từ Person.

Còn Role của bạn khác chính ở đây.

Giả sử:

Manager≺Person,Manager,Person∈Role.

và state có:

σ
Role
	​

(Manager)={m
1
	​

},σ
Role
	​

(Person)={p
1
	​

}.

Ta không muốn:

σ
Att
	​

(email)(m
1
	​

)

tạo ra một storage riêng cho Manager.

Thay vào đó phải có một liên kết:

m
1
	​

⟶p
1
	​

,

và Manager.email được diễn giải thành:

m
1
	​

.email=σ
Att
	​

(email)(p
1
	​

).
	​


Bây giờ giả sử cùng một p
1
	​

 hiện thực hóa thêm một Role khác:

Employee≺Person,

với

σ
Role
	​

(Employee)={e
1
	​

},

và:

m
1
	​

⟶p
1
	​

,e
1
	​

⟶p
1
	​

.

Khi đó:

m
1
	​

.email=e
1
	​

.email=σ
Att
	​

(email)(p
1
	​

).
	​


Đây chính là điều bạn vừa mô tả: hai Role khác nhau nhìn thấy cùng một thuộc tính vì cả hai đang delegate tới cùng một Role instance p
1
	​

.

Nếu qua Manager ta cập nhật:

m
1
	​

.email:=‘‘new@mail.com",

thì về mặt ngữ nghĩa thực sự là:

σ
Att
	​

(email)(p
1
	​

):=‘‘new@mail.com".

Do đó sau đó truy cập qua Employee:

e
1
	​

.email

cũng nhận:

‘‘new@mail.com".

Vậy khác biệt có thể cô đọng thành:

Class inheritance
Role inheritance
	​

child object tự mang gi
a
ˊ
 trị inherited attribute
child role truy cập attribute qua parent-role instance
	​

σ
Att
	​

(a)(x)
σ
Att
	​

(a)(parent
σ
	​

(x))
	​

	​


Hay nói theo cách gần với thiết kế của bạn nhất: Role inheritance ở tầng cú pháp trông giống inheritance, nhưng ở tầng state nó được hiện thực hóa như một association ẩn; việc truy cập qua association đó được làm trong suốt để người dùng vẫn viết r.email như đang truy cập một thuộc tính kế thừa.

Đó cũng là lý do Att
r
∗
	​

 vẫn có thể giữ nguyên như Class:

Att
r
∗
	​

=Att
r
	​

∪
r
′
∈parents(r)
⋃
	​

Att
r
′
	​

,

nhưng σ
Att
	​

 cho inherited attribute của Role phải được diễn giải khác.

thế là viết công thức kế thừa của role vấn giống class đúng không? ta chỉ sửa lại ở phần state thôi phải không? 

Đúng. Ở tầng cú pháp, công thức kế thừa của Role vẫn có thể giữ giống hệt Class:

parents(c)={c
′
∈Class∣c≺c
′
},

và:

Att
c
∗
	​

=Att
c
	​

∪
c
′
∈parents(c)
⋃
	​

Att
c
′
	​

,
navends
∗
(c)=navends(c)∪
c
′
∈parents(c)
⋃
	​

navends(c
′
).

Áp dụng cho mọi:

c∈Class=Entity∪Group∪Role.

Tức là về mặt descriptor, nếu

Manager≺Person,

thì Manager vẫn nhìn thấy toàn bộ attribute và navigable end của Person.

Khác biệt chỉ xuất hiện ở Interpretation/System State:

Entity/Group: inherited attribute được đ
a
ˊ
nh gi
a
ˊ
 trực ti
e
ˆ
ˊ
p tr
e
ˆ
n child instance
	​


còn:

Role: inherited attribute được đ
a
ˊ
nh gi
a
ˊ
 th
o
ˆ
ng qua parent-role instance li
e
ˆ
n k
e
ˆ
ˊ
t với child role
	​


Nói ngắn gọn:

syntax inheritance: gi
o
ˆ
ˊ
ng nhau

nhưng:

state semantics: kh
a
ˊ
c nhau.

Vì vậy phần Generalization hiện tại của bạn không cần tách riêng Role; chỉ khi viết lại System State và σ
Att
	​

 thì mới cần phân nhánh ngữ nghĩa.

Trò chuyện thoại đã kết thúc